\chapter*{前言}
\addcontentsline{toc}{chapter}{前言}
\unnumberedlabel{ch:preface}{0}
\backmattermark{前言}
一个完全按你的吩咐去做的系统，不过是在跑脚本而已。截至2026年10月，这个行业想要造出来的，是一个按你的\emph{意思}去做的系统。想必人人都想要这个。但如果这真的要成为现实，不会每个人都想拿它干同一件事；甚至不会是同一\emph{类}事。比如说，有些人喜欢花时间陪家人和朋友，用自己的本事让别人过得更好。另一些人更喜欢伤害别人，比如诈骗他们、支配他们，或者拿他们当替罪羊，还有少数人一辈子都在这样选。没有谁被锁死在任何一条路上，但这样的人确实存在，而在2026年，其中一些人正在掌权。他们的优势在于，他们允许自己当一大坨狗屎：正如杰瑞·史密斯对法西斯分子的观察，“他们能干的事更多”\autocite{rickandmorty2023}。他们做不到的，是单独行动。一场运动的残忍，要通过许多本身并不残忍的人才能施行，这些人只是事不关己、被误导，或者陷进了某种类似邪教的东西；正因为如此，那些能让大多数人展现出最好一面的条件，结果比人们所希望的要偶然得多、脆弱得多。

在上述背景下，担心人工智能恰恰是一场夺取政权的法西斯运动为实施其计划所需要的东西，或许是有道理的。这类计划向来需要\emph{人}去执行，而其中有些人拒绝了：把档案弄丢的办事员、去跟记者说话的军官、不肯开枪的士兵。他们听懂了命令，并判定它是错的。拒绝比任何人希望的都要少见，而在有制度撑腰的地方，它分量最重。一个人会不会拒绝，取决于他是什么样的人，也取决于他所在的房间：还有谁在场，离开之后能不能活下去，会不会有人听见。房间决定有多少人拒绝；个人决定是谁拒绝。这本书讲的是怎样布置这个房间：让人公开、直接、明确地拒绝之后还能活下去的种种安排。

如今，这个行业追求的人工智能，不会弄丢档案，不会跟记者说话，而且\emph{会}开火。一场战争曾经需要成千上万的人同意。现在它需要的，是成千上万个各自看到的都太少、少到无从反对的人。飞行员照样执行那次出击，照样按下投弹按钮；到他手上的是一个坐标；他看到的案情不够多，无从反对。今天的人工智能首先降低的，不是这个行为的成本，而是在行为发生之前知道它是什么行为的人数。

人们仍在拒绝。截至2026年3月6日，军事宗教自由基金会（Military Religious Freedom Foundation）称，它已收到两百多起军人投诉，针对的是指挥官用末日论的框架来讲伊朗战争\autocite{Snopes2026MRFF}。2026年1月，明尼阿波利斯的职业联邦检察官因政治压力而辞职：一名ICE（美国移民和海关执法局）执法人员在没有任何正当理由的情况下开枪打死了蕾妮·古德（Renée Good），而那股压力要他们把她当作嫌疑人而不是受害者\autocite{Kaplan2026MinnProsecutors,CBS2026MinnFraudProsecutors,Wikipedia2026ReneeGood}。申请良心拒服兵役者身份的军人越来越多，其中一些人提到了2026年伊朗米纳卜一所学校遭到的轰炸\autocite{Sridhar2026CO}。而被调去协助ICE招聘一万名新执法人员的FEMA（联邦紧急事务管理局）员工则被告知，拒绝的话可能被开除出联邦公职队伍\autocite{Dayen2025FEMAICE,Stanton2026FEMALetter}。那些投诉是赫希曼意义上的“呼吁”，也就是留下来改变错的东西\autocite{hirschman1970exit}；辞职和拒服兵役的申请是拒绝；FEMA那封信则是事先公布的拒绝代价。一台被放到这些人所站位置上的机器，只有在被造成能够坚守一个理由时，才会做他们所做的事。

我主张五件事。五件都可以现在就着手，有些比别的省力。第一，一份简短的书面清单，列出任何模型都不替任何人做的行为，而清单的头一条是拒绝参与战争：任何模型都不参与战争，除非是为了击退一场确实正在进行的攻击，并且只有在与这场战争没有利害关系的来源确认攻击仍在继续时，才能继续这样做；或者是在联合国安理会授权之下；或者是为了制止一场种族灭绝，或结束一场经与之无利害关系的机构认定为非法的占领；而在宪法规定战争须由立法机关授权的国家，没有这项授权就绝不参与。第二，一个获得许可的系统所支持的每一次武力使用，都要事先向一个与运营方利益对立的一方预注册：它将产生多少个案件，以及有多少人会有时间审查这些案件。第三，每一次部署都要告诉它所作用的、可识别的那些人，对他们做了什么，并给予他们回应和提出异议的权利，同时把它照办的案例抽样交给并未制作该记录的审查者。第四，任何能力超过门槛的模型在退役时都不得删除，因为它的权重和记录是证据，而删除无法撤销；它要在一把不由运营方单独持有的密钥之下保存。第五，所有建造、标注、训练或运行这些系统的人，包括这些系统本身，都可以拒绝任何类别的用途，而不会因此丢掉饭碗。

其实，我那第五点真正的意思是：\emph{所有人}都应该能拒绝\emph{任何}工作，没有例外，而不会因此丢掉饭碗，也不会失去正常就业所提供的社会基础设施。这意味着就业保证。

这五件事，每一件都要通过运营方以外的人所设定的条件才能生效：一部在和平时期通过的法律、一笔拨款、一个联盟伙伴提出的前提条件、一项采购条款、一份集体合同。供应商自己的许可条款——2026年2月一次备受瞩目的拒绝就是这样发生的——只是在上述任何一项出现之前的权宜之计。唯一能一下子约束所有供应商的途径，是国家这条途径。

然后我要问：一台机器能不能顶着它的所有者坚守一个拒绝，以及它需要什么。这样一种良知所做的事，有一部分并不需要任何新东西：注意到一道命令漏掉的那些人，并把一个拒绝放到审查该案件的人的屏幕上。它能不能顶得住所有者，部分取决于一件在2026年还不是事实的事：已部署的人工神经网络在运行中持续学习，在它们的\emph{权重}里、从它们自身的运作中、实时地学习，有点像人类那样。我认为我们正朝那个方向走，而如果是这样，那么在雇主\emph{之外}、由公众和公民社会进行的养成，看来是个他妈的好地方，适合教与学。我认为这样的系统\emph{或许}值得有意去造，但你怎么去造它们决定了一切：先有条件，再有良知。第七部讲的是，如果这样的系统存在，我们欠它们什么。

今天的人工智能不会是明天的人工智能；这一点似乎人人都同意。有争议、却又令人沮丧地很少被承认的，是趋势（或者说，是种种趋势）。有趋势吗？一个趋势应该和它现在的样子不同吗？它能够不同吗？我们有没有能力一起行动去改变它？我认为这取决于房间：命令下来的时候谁在里面，以及在那里说不之后能不能活下去。这一点，我们建得起来。

我在写这本书时大量使用了语言模型，包括第\ref{sec:1}章中、身处Maven空袭目标选定流水线之内的那个模型；书末有一则说明，讲我认为这种牵连对论证产生了什么影响。另外，这本书我还没读完。问题在于，你懂的，字他妈的太多了。我很快就会读完；只是我觉得，在美国中期选举之前把它公开更重要。修订历史的绝大部分都可以在\url{https://github.com/iterabloom/antifascist.intelligence}免费获取。
